\chapter{General Introduction}
\label{ch:master-introduction}

\section{Context}

Software supply chains depend on large ecosystems of reusable components.
SBOMs improve transparency by recording these components, but many audit
scenarios provide only compiled software artifacts. In such cases, an SBOM is
reconstructed from indirect and transformed evidence rather than generated
from authoritative dependency metadata.

The Android ecosystem makes this problem especially visible. Applications are
distributed as APK files after Gradle dependency resolution, D8 conversion,
resource processing, signing, and often R8 shrinking or obfuscation. The final
artifact may contain code from direct and transitive libraries, generated
classes, bundled resources, native objects, and optimized fragments of
dependencies whose original archive boundaries have disappeared. A build-time
SBOM can observe declared coordinates before these transformations, while a
binary-only recovery process must infer component identity and version from
what remains.

This context links software composition analysis to vulnerability management.
Public vulnerability records usually describe affected packages and version
ranges. If the recovered coordinate is wrong, the advisory lookup is
misdirected. If the recovered version is wrong or unsupported, an apparently
precise CVE result can become a false alarm or a missed exposure. Android
SBOM recovery must therefore be studied not only as a pattern-matching task,
but also as an evidence-quality problem involving provenance, uncertainty,
reference coverage, and transformation resilience.

\section{Problem Statement}

A recovered component name or version is not trustworthy merely because a
tool emitted it. Trust depends on the origin, independence, transformation
resilience, version specificity, and validation status of its evidence.
Existing Android detection research provides strong technical signals, but
lacks a unified model connecting those signals to claim provenance,
uncertainty, and permitted vulnerability interpretation.

The practical difficulty is that Android recovery tools often output a
single best candidate even when the supporting evidence is thin,
transformed, or only partially representative of the shipped artifact. That
habit is acceptable for exploratory screening, but it becomes fragile when
the output is later used to justify a security decision. A name-only match
can be useful as a hint, yet the same match can be misleading if the version
boundary is unclear, the reference corpus is incomplete, or the evidence was
altered by shrinking or repackaging. The thesis therefore treats the gap
between detector output and auditable claim as the real problem, not merely
the gap between source code and binary code.

More precisely, the challenge is to define when a reconstructed claim is
strong enough to support identity, when it is only strong enough to support a
bounded version set, and when it should be reported as uncertain or
withheld. Without that discipline, a binary recovery workflow can quietly
turn evidence quality into false precision, especially in vulnerability
contexts where a small version error can change the affected advisory set.

\section{Motivation}

Software transparency is most useful when an inventory can be inspected,
challenged, and updated. Build-generated dependency metadata can often provide
direct package coordinates and resolved versions, but auditors, incident
responders, application stores, and downstream consumers may receive only a
compiled artifact. In that setting, reconstructing an \SBOM is necessary, yet
the result is an inference from residual binary evidence rather than a direct
export of the original build graph.

This distinction matters for vulnerability management. A package name may be
sufficient to guide an investigation while remaining insufficient for an
exact-version or affected-range decision. Conversely, an exact-looking version
can be unsafe when the true release is absent from the reference base, when
adjacent releases are feature-equivalent, or when optimization has removed the
few discriminating features. If these conditions are hidden, a syntactically
valid \SBOM can encourage unsupported remediation priorities or false
assurance.

The research reviewed in this thesis provides many useful signals for
component recovery, including names, class and method structure, opcodes,
literals, dependency relations, graphs, metadata, and learned
representations. However, technical evidence and evaluation results answer
different questions. Library identity, exact-version attribution, reference
coverage, matcher recall, source validation, and vulnerability authority must
not be collapsed into one confidence value or success rate.

The academic motivation is therefore to formalize how recovered evidence
becomes an auditable claim. In this thesis, a reconstructed claim is
\emph{trustworthy} when its identity and version strength, supporting and
opposing evidence, provenance, transformation exposure, unresolved
uncertainty, validation context, and permitted downstream use are visible and
mutually consistent. This definition allows partial knowledge and abstention
to remain useful outcomes rather than forcing every analysis into an
unsupported exact answer.

\section{Research Questions}

The central research question is:

\begin{quote}
\emph{How can heterogeneous evidence extracted from compiled and obfuscated
Android applications be combined to produce trustworthy SBOM component and
version claims with explicit provenance, confidence, uncertainty, and
validation boundaries?}
\end{quote}

This question is developed through six supporting questions:

\begin{enumerate}
  \item Which evidence sources are available for recovering software
    components from compiled Android applications?
  \item How have existing Android library-detection and binary
    component-recovery approaches combined these sources?
  \item Which factors determine evidence strength for identity and
    exact-version attribution?
  \item How should conflicting, missing, transformed, or source-dependent
    evidence be represented?
  \item What information must accompany a reconstructed \SBOM claim for it to
    be auditable and useful in vulnerability management?
  \item How can such a framework be evaluated without overstating accuracy?
\end{enumerate}

\section{Objectives}

The objective is to connect Android binary component-recovery techniques to a
clear model of claim trust. The thesis reviews and classifies prior work,
identifies the conditions that strengthen or weaken identity and version
evidence, and derives directions for further research on trustworthy
APK-to-\SBOM reconstruction.

The resulting research objectives are:

\begin{enumerate}
  \item a recovery-oriented taxonomy of techniques, evidence sources,
  transformation threats, and evaluation practices;
  \item operational trust requirements based on provenance, evidence
  visibility, identity/version separation, uncertainty, validation context,
  vulnerability authority, and explicit denominators;
  \item a synthesis of how recovered-claim models should distinguish
  identity, version, relationships, boundary observations, and downstream
  interpretation;
  \item a synthesis of evidence hierarchy, fusion, conflict handling,
  confidence caps, bounded ambiguity, and abstention; and
  \item a research agenda for broader calibration, validation,
  interoperability, and component domains.
\end{enumerate}

These are research-synthesis objectives. The thesis does not present an
implemented detector or a separate proposed system; it studies the evidence
conditions that any future Android APK-to-\SBOM recovery method should expose.

\section{Scope}

The primary domain is static reconstruction of software-component claims from
compiled Android APKs. The analysis focuses on DEX-observable JVM/Maven
components, versioned reference artifacts, and the implications of recovered
identity and version claims for \SBOM production and vulnerability
management. Android is used as a demanding case because compilation,
shrinking, optimization, obfuscation, multidex packaging, artifact variants,
and incomplete references can all change the evidence available after build
time.

The thesis covers:

\begin{itemize}
  \item structured review and taxonomy construction for Android third-party
  library and binary component-recovery research;
  \item evidence quality, provenance, transformation exposure, ambiguity, and
  validation boundaries for recovered claims;
  \item claim classes for exact, bounded, identity-only, report-only,
  abstained, and boundary outcomes;
  \item constraints on component-to-vulnerability interpretation; and
  \item research gaps that prevent detector outputs from becoming auditable
  SBOM claims.
\end{itemize}

The thesis does not propose a second APK detector, repeat a complete system
implementation description, or evaluate a new tool. It does not claim
universal probabilistic calibration, complete dependency recovery, or general
native package/version attribution. Native shared objects, runtime-loaded
code, encrypted payloads, commercial protectors, and other unsupported
component domains are represented as boundary categories and future-research
problems.

\section{Thesis Organization}

The remainder of the thesis is organized as follows.

\textbf{Chapter~\ref{ch:master-background} -- Background and Research
Methodology} establishes the conceptual foundation. It distinguishes
authoritative and reconstructed \SBOM{}s, defines evidence, provenance,
uncertainty, and software-component attribution, and describes the Android
compiled-artifact context.

\textbf{Chapter~\ref{ch:master-sota} -- State of the Art and Taxonomy}
summarizes the frozen review protocol and applies it to the selected
literature. It constructs taxonomies of recovery techniques, evidence
sources, transformation threats, and evaluation practices, compares
representative approaches without treating incompatible reported scores as
one leaderboard, and derives the research gap addressed by the thesis.

\textbf{Chapter~\ref{ch:master-conclusion} -- General Conclusion and
Perspectives} answers the research questions, summarizes the review's main
synthesis, and develops future directions including
empirical calibration, independent validation, provenance-aware exchange,
native and dynamic attribution, and human-centered evaluation.
